\section{Introduction and Related Work} \label{sec:introduction}
Linked Open Usable Data stresses that linked data should be consumable by developers~\cite{newbury2018loud}. However, in practice, teams still struggle to bridge ontology engineering and software delivery~\cite{9963753,avila2020applications}. The documentation burden described in NeOn~\cite{suarez2012neon} and later industrial ontology-engineering work~\cite{poveda2022lot} becomes even more challenging when both the ontology, requirements and the codebase evolve together.

Existing tooling only partly covers this gap. Widoco, Parrot, and OnToology generate ontology-centric documentation and evaluation reports~\cite{garijo2017widoco,parrot,alobaid2018ontoology}; other approaches derive specific artefacts such as UML views~\cite{ROBLES201272}, object-oriented libraries~\cite{elkaed2017olga}, REST APIs~\cite{espinoza2022extending}, or SHACL shapes~\cite{cimmino2020astrea}. Broader schema and model-transformation ecosystems, such as the Flemish OSLO toolchain, also support model-to-artefact workflows~\cite{10.1093/gigascience/giaf152,osloUmlTransformer2026,opTedModel2owl2026}. Our solution differs by keeping a governed OWL ontology as the formal source, using SKOS to add project-specific developer names, and projecting one enriched representation to code, SQL, SHACL, diagrams, and documentation in the same regeneration step. The remaining challenge for OWL-first projects is therefore not a single missing generator, but consistently aligning ontology semantics, naming conventions~\cite{7332450,10.1007/978-3-030-38919-2_35}, and implementation constraints across several developer-facing artefacts.

ODDToolkit was built for projects where an ontology must stay aligned with developer-facing artefacts. It addresses three practical needs: a naming layer that projects ontology predicates to implementation names, a shared pipeline that derives multiple artefacts from one enriched model, and integration with normal build workflows so regeneration happens together with ontology changes. The toolkit targets governed ontologies rather than arbitrary third-party vocabularies and assumes project conventions such as SKOS labels, Hydra IRI templates, and explicit enough relation semantics for downstream generation.

This scope also implies a number of structural assumptions. ODDToolkit is aimed at ontologies under active governance, where a project can provide the necessary information to derive cardinality constraints, naming and relation semantics. The toolkit therefore prioritises implementation alignment and regeneration within a controlled development process over full automation for arbitrary third-party vocabularies.